\section{Further finite structural results and local search boundaries}
\label{sec:september30-ledger}

This section records completed results at the fixed intake, not progress
towards a known total number of unrestricted order-$18$ cases. The
finite prerequisites below are supplied as executable certificates.
Historical solver exclusions are distinguished from fresh finite replays.

\begin{corollary}[Homogeneous near-type families]
\label{cor:near-six-bound}
For $n\geq6$, $n\equiv2\pmod4$, a permutation family with all distinct
relative types $\lambda_n$ has at most six members. The bound is sharp
at six. At order eighteen the maximum is exactly four.
\end{corollary}
\begin{proof}
Normalize two members to $\mathrm{id},A$ and fix the four-point support
$S$ of $A$'s four-cycle. By Theorem~\ref{thm:near-core}, any other
member preserves a six-point component containing $S$, so sends at
most two of its points outside $S$. For a family of size $k$, its
images at each point are distinct. Consequently
$4(k-4)\leq2(k-2)$, hence $k\leq6$.
The following Latin6 table attains six:
\[
\begin{pmatrix}
0&1&2&3&4&5\\1&2&3&0&5&4\\2&0&4&5&1&3\\
3&4&5&1&0&2\\4&5&0&2&3&1\\5&3&1&4&2&0
\end{pmatrix}.
\]
This control has pattern FTF, not FFF.

At eighteen the complete third-line domain has
$7\cdot16\cdot15\cdot2^3=13440$ members: choose the transposition
pair which joins $S$, a six-point template, a pairing of the remaining
six transposition pairs, and their two transverse orientations.
Root-centralizing conjugations partition it into two orbits of size
$6720$. For each representative exactly $39$ of all $13440$ possible
fourth lines are compatible, and each has the same exceptional core.
The independent physical-matching checker verifies this complete
finite common-core premise, not just sample neighbors.

If $k\geq5$, every member therefore preserves one six-point core
and its twelve-point complement. Each normalized nonidentity core
restriction has type $4+2$, hence positive sign. Every relative core
restriction is fixed-point-free, has only even-length cycles, and has
positive sign; therefore its type is $4+2$. On the complement all relatives are
fixed-point-free involutions. The normalized restrictions consequently
commute and generate an elementary abelian $2$-group. Each orbit
contains the $k$ distinct images of any one point, so its power-of-two
size is at least eight. Twelve cannot be a sum of such orbit sizes.
Thus $k\leq4$. For sharpness take the first four rows of the displayed
table and adjoin three four-point blocks on which row $i$ acts as
$x\mapsto x\mathbin{\mathrm{xor}}i$, $0\leq i<4$.
This is a Latin rectangle, not an FFF18 table.
\end{proof}

The four-line bound applies to any Latin square in each physical view.
Under full FFF, Theorem~\ref{thm:near-triangle} is stronger: already
three pairwise near lines are impossible. An exact one-view diamond
spectrum is retained in the certificate, but its five-near-edge
antecedent is therefore not a surviving FFF search branch.

\begin{computationaltheorem}[A conditional mixed-core family]
\label{thm:mixed-core}
Let $n=6+4t$, $t\geq0$, and fix a fixed-point-free involution $A$.
The permutations $B$ with types
\[
 B:4+2^{(n-4)/2},\qquad A^{-1}B:6+2^{(n-6)/2}
\]
have one six-point component and $t$ four-point components. They form
one orbit under the centralizer of $A$, with
\[
 \binom{3+2t}{3}\,24(2t-1)!!\,2^t
\]
members and stabilizer order $2\cdot4^t t!$, where $(-1)!!=1$.
Their three-row rectangles have no already-closed odd companion cycle.
\end{computationaltheorem}
\begin{proof}
Use the three-matching component inequality from
Theorem~\ref{thm:near-core}, now counting deficit from involutions.
The four-cycle contributes one and the six-cycle two. They cannot
occur in different components: the product of two fixed-point-free
involutions has paired equal-length cycles and cannot contain the
isolated exceptional cycle. For a component of size $d$ and total
cycle deficit $\delta$, the inequality and sign identity give
\[
d\leq2\delta+4,\qquad \delta\equiv d/2\pmod2,\qquad d\geq4\text{ even}.
\]
Thus the component with deficit three has size six or ten; components
with deficit zero have size four. A connected ten-point component,
if present, would restrict to a candidate in the complete enumeration below.

At degree ten a complete forward enumeration of $18900$ near
permutations and a separate reverse enumeration of $75600$
single-six-cycle permutations give exactly the same $480$ candidates.
All have component sizes $6+4$. At degree six there are exactly
$24$ candidates, one orbit of the $48$-element matching centralizer.
These finite premises are explicitly replayed in the package.

Choose three of $A$'s $3+2t$ pairs for the exceptional core, one
of the $24$ local templates, and a pairing of the other $2t$ pairs
into four-point blocks, each with two orientations. This is unique
and exhaustive. The centralizer transports all choices to one another;
its local stabilizers have orders two and four, with the $t$ ordinary
blocks freely permutable. This proves the count and stabilizer.

All $24$ six-point cores and both ordinary four-point blocks have
no closed odd column or symbol cycle, as checked by direct physical
matchings. For a pair across two distinct components, each specified
arrow runs from one disjoint component to the other and cannot close.
Thus their union has no already-closed odd companion cycle.
\end{proof}

At eighteen the count is $241920$ and the stabilizer order is $768$.
This gives a canonical prefix only \emph{if} a square contains the
specified triple; it does not force its occurrence or prove completion.
Its involutive second row is negative and cannot be combined with the
positive-anchor normalizations in Corollary~\ref{cor:nonnear-anchors}.
A retained clean $12\times18$ rectangle has a proof-checked exclusion
of extension to thirteen rows with all twelve rows fixed. This is a
dead leaf, not an exclusion of the three-row family. Relaxed and full
completion pilots time out. No clean thirteen-row extension or FFF18
square is established by this intake.

\subsection{Reusable first-touch argument and numerical applications}
\begin{lemma}[First-touch support bound]
Fix a labelled Latin table and a collection of its odd-cycle cell
supports. During a sequence of two-line cycle-union trades, a support
which is touched for the first time must meet one of the two selected
view-tagged line labels $(\text{view},\text{label})$, evaluated on the original
support. If every possible set of the at most $2k$ selected line
labels leaves an original support uncovered, $k$ such trades cannot
reach FFF.
\end{lemma}
\begin{proof}
Before its first touch every cell and entry of the support is unchanged.
A row or column trade affects only its two chosen rows or columns.
A symbol trade affects only occurrences of its two selected symbols;
on the untouched support these are still the original symbols.
Therefore the first touch is incident with one of the corresponding
original labels. A support missed by every selected label is never
touched, so all arrows of its original odd cycle persist. This argument
allows adaptive trades in mixed views and overlapping changed regions.
\end{proof}

For one frozen FTF18 source, thirty-six pairwise disjoint column
three-cycle supports are retained. Each row or symbol label meets
six, and each column label four. Exhausting all $\binom{36}{6}=1947792$
six-label row/symbol selections gives maximum coverage $33$, by two
independent incidence calculations. Three moves including a column
move cover at most $4\cdot6+2\cdot4=32$. Hence at most three moves
cannot make even its column view F: distance is \emph{at least four},
not exactly four. A static four-move support cover is not a valid
trade sequence or a construction.

Two further retained applications concern precisely the $12513$ odd
translation-equivariant $\mathbb F_{17}$ complete-mapping one-point
tables, not all order-$18$ squares. Their unchanged-line vertex-cover
bounds exclude up to six intercalate trades from every source, and
up to seven from $9537$ of them. A separate one-step argument excludes
every proper two-line cycle-union trade: the residual affine F-view
domain has $8160$ moves and $7599$ distinct labelled outputs, none FFF.
The reported minimum of $64$ failing pairs pertains only to those
enumerated outputs, not to every source ruled out by the reduction.
These are bounded negative applications; neither more general moves
nor other starting families are excluded.

\paragraph{Verification and open complement.}
The new palette, mixed-core and lift results have fresh finite replays.
The retained source-specific trade-radius computations have fresh
rechecks in the private research checkpoint; the public numerical
ledger labels their original source and coverage rather than presenting
them as new global theorems. The historical DRAT leaf above is not
freshly recertified by the portable paper checker. Timeouts, necessary
conditions, and complete counts without mate coverage do not decide
unrestricted order $18$. Literature priority and independent expert
review remain separate, uncompleted tasks.
